\SourcePage{118}{121}
\section{Angular-Momentum Eigenfunctions}

Specifying the values of $l$ and $m$ does not determine a particle's wave
function completely. Indeed, the expressions for the operators of these
quantities in spherical coordinates contain only the angles $\theta$ and
$\varphi$, so their eigenfunctions may include an arbitrary factor depending
on $r$. Here we shall consider only the angular part of the wave function
characteristic of angular-momentum eigenfunctions. Denote it by
$Y_{lm}(\theta,\varphi)$ and impose the normalization condition
\[
  \int |Y_{lm}|^2\,do=1
\]
($do=\sin\theta\,d\theta\,d\varphi$ is the element of solid angle).

The calculation below shows that the problem of finding the common
eigenfunctions of $\hat{\mathbf l}^2$ and $\hat l_z$ permits separation of
the variables $\theta$ and $\varphi$. These functions may therefore be
sought in the form
\begin{equation}
  Y_{lm}=\Phi_m(\varphi)\Theta_{lm}(\theta),
  \tag{28.1}
\end{equation}
where $\Phi_m(\varphi)$ are the eigenfunctions of $\hat l_z$ defined by
(27.3). Since the functions $\Phi_m$ have already been normalized by (27.4),
the functions $\Theta_{lm}$ must satisfy
\begin{equation}
  \int_0^\pi |\Theta_{lm}|^2\sin\theta\,d\theta=1.
  \tag{28.2}
\end{equation}
Functions $Y_{lm}$ with different values of $l$ or $m$ are automatically
mutually orthogonal:
\begin{equation}
  \int_0^{2\pi}\!\int_0^\pi
  Y_{l'm'}^*Y_{lm}\sin\theta\,d\theta\,d\varphi
  =\delta_{ll'}\delta_{mm'},
  \tag{28.3}
\end{equation}
\SourcePage{119}{122}
because they are eigenfunctions of angular-momentum operators belonging to
different eigenvalues. The functions $\Phi_m(\varphi)$ are also separately
orthogonal (see (27.4)), since they are eigenfunctions of $\hat l_z$
belonging to its different eigenvalues $m$. The functions
$\Theta_{lm}(\theta)$ by themselves, however, are not eigenfunctions of any
angular-momentum operator; they are mutually orthogonal for different $l$,
but not for different $m$.

The most direct way to calculate the required functions is to solve the
eigenfunction problem for $\hat{\mathbf l}^2$ written in spherical
coordinates (formula (26.16)). The equation
$\hat{\mathbf l}^2\psi=l^2\psi$ reads
\[
 \frac1{\sin\theta}\frac\partial{\partial\theta}
 \left(\sin\theta\frac{\partial\psi}{\partial\theta}\right)
 +\frac1{\sin^2\theta}\frac{\partial^2\psi}{\partial\varphi^2}
 +l(l+1)\psi=0.
\]
Substitution of $\psi$ in the form (28.1) gives the following equation for
$\Theta_{lm}$:
\begin{equation}
 \frac1{\sin\theta}\frac d{d\theta}
 \left(\sin\theta\frac{d\Theta_{lm}}{d\theta}\right)
 -\frac{m^2}{\sin^2\theta}\Theta_{lm}
 +l(l+1)\Theta_{lm}=0.
 \tag{28.4}
\end{equation}
This equation is well known from the theory of spherical harmonics. It has
solutions satisfying the finiteness and single-valuedness conditions for
nonnegative integral values $l\geq|m|$, in agreement with the
angular-momentum eigenvalues obtained above by the matrix method. The
corresponding solutions are the associated Legendre polynomials
$P_l^m(\cos\theta)$ (see \S~c of the mathematical supplements). Normalizing
the solution by (28.2), we obtain\footnote{The normalization condition does
not, of course, determine the choice of phase factor. The convention used in
this book is the one most natural from the standpoint of the general theory
of angular-momentum addition; it differs from the convention usually adopted
by a factor $i^l$. The advantages of this choice will be apparent from the
notes on pp.~278, 527, and 533.}
\begin{equation}
 \Theta_{lm}=(-1)^mi^l
 \sqrt{\frac{(2l+1)(l-m)!}{2(l+m)!}}P_l^m(\cos\theta).
 \tag{28.5}
\end{equation}
Here $m\geq0$ is assumed. For negative $m$, define $\Theta_{lm}$ by
\begin{equation}
  \Theta_{l,-|m|}=(-1)^m\Theta_{l|m|},
  \tag{28.6}
\end{equation}
that is, for $m<0$, $\Theta_{lm}$ is given by (28.5) with $|m|$ in place of
$m$ and with the factor $(-1)^m$ omitted.
\SourcePage{120}{123}
Thus, from the mathematical point of view, the angular-momentum
eigenfunctions are spherical harmonics with a particular normalization. For
later reference, we write their complete expression, incorporating all the
definitions just given:
\begin{equation}
 Y_{lm}(\theta,\varphi)=(-1)^{(m+|m|)/2}i^l
 \left[\frac{2l+1}{4\pi}\frac{(l-|m|)!}{(l+|m|)!}\right]^{1/2}
 P_l^{|m|}(\cos\theta)e^{im\varphi}.
 \tag{28.7}
\end{equation}
In particular,
\begin{equation}
  Y_{l0}=i^l\sqrt{\frac{2l+1}{4\pi}}P_l(\cos\theta).
  \tag{28.8}
\end{equation}
Functions differing only in the sign of $m$ are evidently related by
\begin{equation}
  (-1)^{l-m}Y_{l,-m}=Y_{lm}^*.
  \tag{28.9}
\end{equation}

For $l=0$ (and consequently $m=0$), the spherical harmonic reduces to a
constant. In other words, the wave functions of particle states with zero
angular momentum depend only on $r$ and hence possess complete spherical
symmetry, in accordance with the general statement made in \secref{27}.

For a fixed $m$, the values of $l$ beginning with $|m|$ enumerate, in
ascending order, the successive eigenvalues of $\mathbf l^2$. The general
theorem on zeros of eigenfunctions (\secref{21}) therefore implies that
$\Theta_{lm}$ vanishes at $l-|m|$ different values of $\theta$; in other
words, its nodal lines are $l-|m|$ “circles of latitude” on the sphere.
For the complete angular functions, if real factors $\cos m\varphi$ or
$\sin m\varphi$ are chosen instead of
$e^{\pm i|m|\varphi}$\footnote{Each such function represents a state in
which $l_z$ has no definite value, but may take the values $\pm m$ with equal
probability.}, they also have $|m|$ “meridian circles” as nodal lines. The
total number of nodal lines is therefore $l$.

Finally, let us show how the functions $\Theta_{lm}$ can be calculated by the
matrix method. The procedure is analogous to the calculation of the
oscillator wave functions in \secref{23}. We start from (27.8),
$\hat l_+Y_{ll}=0$. Using (26.15) for the operator $\hat l_+$ and substituting
\[
  Y_{ll}=\frac1{\sqrt{2\pi}}e^{il\varphi}\Theta_{ll}(\theta),
\]
\SourcePage{121}{124}
we obtain for $\Theta_{ll}$ the equation
\[
  \frac{d\Theta_{ll}}{d\theta}-l\cot\theta\cdot\Theta_{ll}=0,
\]
whence
$\Theta_{ll}=\operatorname{const}\cdot\sin^l\theta$. Determining the constant
from the normalization condition gives
\begin{equation}
 \Theta_{ll}=(-i)^l\sqrt{\frac{(2l+1)!}{2}}\frac1{2^ll!}\sin^l\theta.
 \tag{28.10}
\end{equation}
Next, using (27.12), write
\[
 \hat l_-Y_{l,m+1}=(l_-)_{m,m+1}Y_{lm}
 =\sqrt{(l-m)(l+m+1)}Y_{lm}.
\]
Repeated application of this formula gives
\[
 \sqrt{\frac{(l-m)!}{(l+m)!}}Y_{lm}
 =\frac1{\sqrt{(2l)!}}\hat l_-^{l-m}Y_{ll}.
\]
The right-hand side is readily evaluated with the aid of (26.15) for
$\hat l_-$, according to which
\[
 \hat l_-\left[f(\theta)e^{im\varphi}\right]
 =e^{i(m-1)\varphi}\sin^{1-m}\theta\frac d{d\cos\theta}
 \left(f\sin^m\theta\right).
\]
Repeated application of this formula gives
\[
 \hat l_-^{l-m}e^{il\varphi}\Theta_{ll}
 =e^{im\varphi}\sin^{-m}\theta
 \frac{d^{l-m}}{(d\cos\theta)^{l-m}}
 \left(\sin^l\theta\cdot\Theta_{ll}\right).
\]
Finally, using these relations together with (28.10) for $\Theta_{ll}$, we
obtain
\begin{equation}
 \Theta_{lm}(\theta)=(-i)^l
 \sqrt{\frac{2l+1}{2}\frac{(l+m)!}{(l-m)!}}
 \frac1{2^ll!\sin^m\theta}
 \frac{d^{l-m}}{(d\cos\theta)^{l-m}}\sin^{2l}\theta,
 \tag{28.11}
\end{equation}
which agrees with (28.5).
